\documentclass[11pt]{article}
% Shared preamble for the Cupcast papers.
\usepackage[a4paper,margin=2.8cm]{geometry}
\usepackage{amsmath,amssymb,mathtools}
\usepackage{booktabs}
\usepackage{graphicx}
\usepackage{microtype}
% Classic BibTeX via natbib: tectonic runs bibtex natively, while its bundled
% biblatex is version-locked against system biber. Each paper folder symlinks
% references.bib from shared/.
\usepackage[numbers,sort&compress]{natbib}
\usepackage{xcolor}
\usepackage[colorlinks=true,allcolors=blue!45!black]{hyperref}
\graphicspath{{../figures/}}

% Shared notation
\newcommand{\E}{\mathbb{E}}
\newcommand{\Prob}{\mathbb{P}}
\newcommand{\att}{\alpha}
\newcommand{\dfn}{\beta}


\title{Cupcast I: A Dixon--Coles Model with Squad-Quality Shrinkage\\for Forecasting the 2026 FIFA World Cup}
\author{Carlos Garcia}
\date{June 2026}

\begin{document}
\maketitle

\begin{abstract}
I specify a forecasting model for the 2026 FIFA World Cup: a Dixon--Coles goal
model \cite{dixoncoles1997} fitted by weighted maximum likelihood on a decade of
international results, whose team ratings are shrunk toward an empirical prior
built from Elo ratings and minutes-weighted, performance-based squad-quality
composites. Match probabilities feed a Monte Carlo simulation of the full
48-team tournament under FIFA's official group tiebreakers and Round-of-32
bracket allocation. All tuning decisions (time-decay rate, friendly weighting,
shrinkage strength) are made by out-of-sample predictive likelihood and
reported here. Validation is the subject of Paper~II; the forecast itself of
Paper~III.
\end{abstract}

\section{Introduction}
The design goals are: (i) every match receives a full scoreline distribution,
not a point prediction; (ii) every modeling choice is either standard in the
literature or selected by out-of-sample evidence; (iii) squad quality enters
through measurable on-pitch output, never market valuations; and (iv) the whole
pipeline is reproducible from a clean checkout --- data fetchers are cache-first,
every stochastic step is explicitly seeded, and the published numbers regenerate
from one command.

The pipeline is: match history $\rightarrow$ Elo ratings (\S\ref{sec:elo})
$\rightarrow$ Dixon--Coles fit (\S\ref{sec:dc}) $\rightarrow$ shrinkage toward a
squad-informed prior (\S\ref{sec:shrink}) $\rightarrow$ 50{,}000 tournament
simulations (\S\ref{sec:sim}).

\section{Data}\label{sec:data}
\subsection{International match history}
The training table is built from the community-maintained CC0 dataset of
international results \cite{martj42}, current through June 9, 2026, totalling
49{,}400 matches since 1872 with venues and neutral-site flags. Cross-checks
and competition metadata (rounds, fixtures with team identifiers) come from
API-Football \cite{apifootball} for the 2022--2024 seasons. Non-FIFA
internationals are removed by restricting to teams that appear in World Cup
qualification or a Nations League within the window; the Dixon--Coles training
set since January 2015 holds 10{,}577 matches across 226 teams.

\subsection{Squads and player data}
Confirmed 26-man squads for all 48 teams (registration deadline June 2, 2026)
are parsed from the structured player templates of the consolidated squads
page \cite{wikisquads}: 1{,}247 players with shirt number, position, age at
kickoff, caps, and club. (Argentina registered 25 players; the source is
faithful on this point.) Player-level performance comes from FBref's
Opta-sourced 2025--26 season tables \cite{fbref} --- minutes, npxG, xAG,
progressive actions, and goalkeeper post-shot expected goals --- covering the
big five European leagues. Coverage tiers and the fallback for squads based in
uncovered leagues are described in \S\ref{sec:squad}. No market-value data is
used anywhere.

\subsection{Bookmaker odds and club data}
Outright championship odds across bookmakers are snapshotted from The Odds API
\cite{oddsapi}, with proportional overround removal; the model--market
comparison appears in Paper~III. Club league results with Pinnacle closing odds
\cite{footballdata} support the engine-level validation tier in Paper~II.

\section{Elo Ratings}\label{sec:elo}
The rating input follows the World Football Elo system \cite{eloratings}, whose
predictive value for match outcomes is established empirically by
\citet{hvattumarntzen2010}. After each match the rating of team $i$ updates as
\begin{equation}
  R_i' = R_i + K \, G \,(W - W_e), \qquad
  W_e = \frac{1}{1 + 10^{-(R_i - R_j + H)/400}},
\end{equation}
where $W$ is the realised result, $H{=}100$ the home adjustment (zero at
neutral venues), $G$ the goal-difference multiplier ($1$, $1.5$, then
$1.75 + (n-3)/8$ for margins $n \ge 3$), and $K$ encodes importance: 60 for
World Cup finals matches, 50 for continental finals tournaments, 40 for World
Cup and continental qualifiers and the Nations Leagues, 30 for other
tournaments, 20 for friendlies. Ratings are computed from the raw match
history rather than scraped, which yields reproducible as-of-date ratings for
backtesting; the June 2026 ladder (Spain 2219, Argentina 2190, France 2125,
England 2086, Brazil 2069, \ldots) matches the published ordering.

\section{The Goal Model}\label{sec:dc}
Goals scored by teams $i$ and $j$ are conditionally Poisson with rates
\begin{equation}
  \lambda_i = \exp\!\bigl(\mu + \att_i - \dfn_j + \gamma\, h_i\bigr), \qquad
  \lambda_j = \exp\!\bigl(\mu + \att_j - \dfn_i + \gamma\, h_j\bigr),
\end{equation}
where $\att$ and $\dfn$ are attack and defense ratings and $h$ is a host
indicator: in training it marks genuine home fixtures; at the 2026 World Cup it
is nonzero only for Mexico, the United States, and Canada playing inside their
own country (the schedule keeps all three at home for the entire group stage).
Low-score dependence is corrected by the Dixon--Coles adjustment
\cite{dixoncoles1997}
\begin{equation}
  \Prob(X{=}x, Y{=}y) = \tau_{\rho}(x,y)\,
    \frac{e^{-\lambda_i}\lambda_i^{x}}{x!}\,
    \frac{e^{-\lambda_j}\lambda_j^{y}}{y!},
\end{equation}
with $\tau_\rho$ reweighting the $\{0,1\}^2$ scorelines, addressing the draw
deficit of the independent-Poisson model of \citet{maher1982}.

\subsection{Estimation}
Parameters are estimated by weighted maximum likelihood with analytic
gradients under L-BFGS-B; identifiability is imposed through a quadratic
penalty pinning the means of $\att$ and $\dfn$ at zero, and a weak ridge
($10^{-2}$) keeps teams with negligible decayed sample bounded. The 483-parameter
fit on 10{,}577 matches converges in roughly one second. Fitted globals as of
June 11, 2026: $\mu = 0.062$, home advantage $\gamma = 0.242$ (a $+27\%$ goal
rate), $\rho = -0.046$.

\subsection{Recency weighting}\label{sec:decay}
Matches are weighted by $\phi(t) = \exp(-\xi\, \Delta t)$. The decay rate and a
candidate friendly down-weight were tuned jointly by out-of-sample log-loss
over eight half-year rolling folds (January 2022 -- June 2026, $n = 4167$
matches):

\begin{center}
\begin{tabular}{lccc}
\toprule
Half-life & friendly $\times 0.3$ & friendly $\times 0.6$ & friendly $\times 1.0$ \\
\midrule
1.5y & 0.9246 & 0.8962 & 0.8737 \\
2.5y & 0.8646 & 0.8631 & \textbf{0.8623} \\
3.5y & 0.8854 & 0.8633 & 0.8626 \\
5.0y & 0.8837 & 0.8643 & 0.8637 \\
\bottomrule
\end{tabular}
\end{center}

The optimum is flat between half-lives of 2.5 and 3.5 years, consistent with
\citet{leyetal2019}; 2.5 years is used. Notably, down-weighting friendlies
\emph{strictly worsens} predictive likelihood at every half-life. The Elo
$K$-hierarchy intuition --- friendlies matter less --- concerns rating
\emph{stakes}, not informational content: friendly scorelines are evidently as
informative about team strength as competitive ones, so they enter the
likelihood at full weight.

\section{Squad-Quality Composite}\label{sec:squad}
Each player's quality is the league-pooled $z$-score of npxG+xAG per 90
minutes over the 2025--26 club season (minimum 450 minutes). Squad composites
aggregate player quality using \emph{expected tournament minutes} as weights.
Expected minutes follow a depth-chart heuristic on a 4-3-3 baseline: within
each position group, starters are the most-capped players; goalkeepers play
the full match, outfield starters retain 85\% of their slot's minutes, and the
bench splits the remainder in proportion to caps. Squad-level age is also
computed minutes-weighted. Players not matched to the covered leagues are
imputed at the matched pool's 20th percentile; a team's \emph{coverage} is the
share of its expected minutes that matched. Teams under 35\% coverage fall
back to the Elo-only prior tier below, and per-team coverage is reported with
the forecast. The caps-based depth chart is a deliberate v1 simplification ---
it under-ranks recently emerged starters --- and is the first refinement
candidate listed in Paper~III's limitations.

\section{Shrinkage Toward an Empirical Prior}\label{sec:shrink}
Fitted ratings are blended with a cross-team empirical prior,
\begin{equation}
  \tilde{\att}_i = w_i\, \hat{\att}_i + (1 - w_i)\, \att_i^{\mathrm{prior}},
  \qquad
  w_i = \frac{n_i^{\mathrm{eff}}}{n_i^{\mathrm{eff}} + k},
\end{equation}
where $n_i^{\mathrm{eff}}$ is the decay-weighted match count. The prior is the
prediction of a weighted least-squares regression of fitted ratings on
standardised Elo --- plus the squad composite where coverage permits --- so that
well-identified teams define the relationship and thinly observed teams are
pulled toward it: the partial-pooling logic of \citet{baioblangiardo2010} in a
transparent two-stage form. The strength $k$ was tuned on the same rolling
folds as \S\ref{sec:decay}:

\begin{center}
\begin{tabular}{lccccc}
\toprule
$k$ & 0 & \textbf{10} & 25 & 50 & 100 \\
Log-loss & 0.8623 & \textbf{0.8615} & 0.8641 & 0.8671 & 0.8706 \\
\bottomrule
\end{tabular}
\end{center}

Light shrinkage ($k{=}10$) helps; anything stronger is worse than none. Under
the tuned decay an international regular accumulates $n^{\mathrm{eff}} \approx 40$,
so established teams keep $\ge 80\%$ of their fitted rating while a thin-history
qualifier leans on the prior.

\section{Tournament Simulation}\label{sec:sim}
Each of the 50{,}000 simulations plays all 72 group matches by sampling
scorelines from the Dixon--Coles grid, ranks groups under the FIFA tiebreaker
chain (points, goal difference, goals for, then head-to-head points/GD/GF among
the tied teams, then drawing of lots --- fair-play points are not modelled, a
documented simplification), ranks the twelve third-placed teams by the same
criteria to find the eight qualifiers, and maps qualifiers into the official
Round-of-32 bracket \cite{wikiknockout}.

\paragraph{Third-place allocation.} FIFA's Annex C \cite{fifaregs} fixes the
assignment of qualified thirds to the eight bracket slots for each of the 495
possible combinations. The published table is not machine-readable; the
simulator instead solves each combination by deterministic constraint matching
(most-constrained slot first, alphabetical preference). This reproduces FIFA's
published worked example exactly, and a unit test verifies that all 495
combinations yield valid assignments. Any divergence from unpublished rows of
Annex C would only permute which group winner faces which third --- a
second-order effect on team-level probabilities.

\paragraph{Knockouts.} Ties after regulation continue into extra time, modelled
as a Dixon--Coles grid at one-third rates, then a penalty shootout
\begin{equation}
  p = \tfrac12 + \operatorname{clip}\bigl(0.02\,(z_a - z_b),\, \pm 0.05\bigr),
\end{equation}
where $z$ is the first-choice goalkeeper's shot-stopping over-performance
(PSxG$\pm$ per 90, standardised across tournament keepers \cite{fbref});
without keeper data the shootout is an even coin. Host advantage applies in
knockouts only when a host plays inside its own country, which the bracket
fixes per match.

\section{Reproducibility}
Every fetcher caches raw responses on first call and the pipeline reads only
from cache; API-Football calls pass through a persisted request ledger. All
random draws flow from explicitly passed seeds (NumPy \texttt{Generator}); the
full pipeline --- fit, shrinkage, 50{,}000 simulations, reports --- runs in well
under a minute on commodity hardware via \texttt{python -m cupcast.report.build},
and the tuning grids in \S\ref{sec:decay} and \S\ref{sec:shrink} regenerate via
the documented validation modules.

\bibliographystyle{abbrvnat}
\bibliography{references}
\end{document}
